% ---
% title: Conditional probability: special cases
% tags: conditional-probability
% ---

\begin{problem}
	Let $A$ and $B$ be events with $P(B)>0$. Determine $P(A\mid B)$ if
	(a) $B$ implies $A$, that is, $B\subseteq A$;
	(b) $A$ implies $B$, that is, $A\subseteq B$;
	(c) $A$ and $B$ are mutually exclusive.
	In each case, also give an example with one roll of a fair die and check
	the result by counting.
\end{problem}

\begin{solution}
	\idea Start from the definition
	\[ P(A\mid B)=\frac{P(AB)}{P(B)}. \]
	In each special case the intersection $AB$ is something simple: $B$, $A$,
	or the empty set.
	\step{(a)} If $B\subseteq A$, then every outcome of $B$ is also in $A$. So
	the outcomes that are in both $A$ and $B$ are exactly the outcomes of $B$:
	$AB=B$. Therefore
	\[ P(A\mid B)=\frac{P(AB)}{P(B)}=\frac{P(B)}{P(B)}=1. \]
	In words: if we know that $B$ happened, then $A$ surely happened too.

	Example: roll a fair die, let $A=\{2,4,6\}$ (``even'') and $B=\{6\}$
	(``six''). A six is always even, so $B\subseteq A$. Given that the die shows
	a six, the result is certainly even: $P(A\mid B)=1$.
	Check by counting: if we know the die shows a six, there is $1$ equally
	likely outcome left ($6$), and it is even. So $P(A\mid B)=1/1=1$.
	\step{(b)} If $A\subseteq B$, then every outcome of $A$ is also in $B$, so
	$AB=A$. Therefore
	\[ P(A\mid B)=\frac{P(AB)}{P(B)}=\frac{P(A)}{P(B)}. \]
	Since $P(B)\le1$, this is at least $P(A)$: learning that $B$ happened never
	makes $A$ less likely. It strictly increases the probability of $A$ when
	$P(A)>0$ and $P(B)<1$.

	Example: $A=\{6\}$ and $B=\{2,4,6\}$. Now $A\subseteq B$, and
	\[ P(A\mid B)=\frac{1/6}{1/2}=\frac13. \]
	Check by counting: if we know the result is even, there are $3$ equally
	likely outcomes left ($2$, $4$, $6$), and exactly one of them is a six.
	So the probability is $1/3$, twice the unconditional $P(A)=1/6$.
	\step{(c)} If $A$ and $B$ are mutually exclusive, they have no common
	outcome: $AB=\emptyset$, so $P(AB)=0$. Therefore
	\[ P(A\mid B)=\frac{0}{P(B)}=0. \]
	In words: if $B$ happened, then $A$ cannot have happened.

	Example: $A=\{1\}$ and $B=\{2,4,6\}$. A one is not even, so the two events
	are mutually exclusive. Given that the result is even, it is certainly not
	a one: $P(A\mid B)=0$.
	Check by counting: if we know the result is even, there are $3$ equally
	likely outcomes left ($2$, $4$, $6$), and none of them is a one. So
	$P(A\mid B)=0/3=0$.
	\begin{answer}
		(a) $1$ \quad (b) $\dfrac{P(A)}{P(B)}$ \quad (c) $0$
	\end{answer}
	\begin{remark}
		\textbf{Common mistake:} mixing up $P(A\mid B)$ and $P(B\mid A)$. In
		the example of (a), $P(A\mid B)=1$, but
		\[ P(B\mid A)=\frac{1/6}{1/2}=\frac13. \]
		A six is always even, but an even result is not always a six. Cases (a)
		and (b) are the same situation seen from the two sides.

		Also, case (c) shows that mutually exclusive events (with positive
		probabilities) strongly influence each other: knowing that $B$
		happened changes the probability of $A$ all the way to $0$.
	\end{remark}
\end{solution}
